\documentclass{enset}
\usepackage[brazil]{babel}      % para texto em Português
%\usepackage[latin1]{inputenc}   % para acentuação em Português OU
\usepackage[utf8]{inputenc}   % para acentuação em Português com o uso do Unicode, 
% mude a codificação desse padrão para utf-8

%%
%% POR FAVOR, NÃO FAÇA MUDANÇAS NESSE PADRÃO QUE ACARRETEM  EM
%% ALTERAÇÃO NA FORMATAÇÃO FINAL DO TEXTO
%%

\title{Instruções para Submissão de Trabalhos para o ENSET 2019}

% Acrescente quantos autores precisar.
% Use letras sequenciais (a,b,c,d...) entre parenteses para identificar as afiliações e use o comando \thanks  para os e-mails, separados por ponto e virgula.
\author[ (a, b)]{Primeira Autora~\thanks{primeira.autora@unital.edu}}
\author[ (b, c)]{Segundo Autor~\thanks{segundo.autor@uninova.edu;  alternativo@gmail.com}}
\author[ (a, c, d)]{Outra Autora~\thanks{outro.autor@uniotima.com.br}}
\author[ (b)]{Mais U.M. Autor~\thanks{maisum@yahoo.com.br}}

% Acrescente quantos autores precisar.
% Use letras sequenciais (a,b,c,d...) para identificar as afiliações.
\affil[a ]{Departamento Exemplo, Universidade Tal, Endereço da Universidade Tal}
\affil[b ]{Departamento Modelo, Universidade Nova, Endereço da Universidade Nova}
\affil[c ]{Centro de Pesquisa, Universidade Ótima, Endereço da Universidade Ótima}
\affil[d ]{Outro Centro de Pesquisa, Universidade Ótima, Endereço da Universidade Ótima}

\begin{document}

\criartitulo

\section{Introdução}

Este é o modelo para submissão de trabalhos para apresentação em formato de painel no ENSET 2019. Poderão ser submetidos trabalhos de pesquisa e extensão, em nível de graduação e pós-graduação, em engenharia, tecnonogia e áreas afins. Os trabalhos submetidos poderão estar em qualquer etapa do se desenvolvimento, isto é, poderão ser apresentados resultados preliminares ou finais. 

A submissão deverá ser realizada na forma de resumo expandido, em Português, que não deve exceder três páginas, incluindo as referências, e deve ressaltar os objetivos, metodologia, resultados e conclusões. Cada um desses tópicos poderão ser divididos em parágrafos, formando um texto corrido, ou divididos em seções, como no presente modelo.

Os trabalhos submetidos que não respeitem a formatação final do texto produzida por este padrão serão \textbf{rejeitados} pelo Comitê Editorial do evento. Para que os trabalhos submetidos sejam publicados nos Anais do Evento é necessário que o autor-apresentador tenha pago a taxa de inscrição até a data definida no crograma do evento. 

Solicitamos que não sejam alteradas as configurações de tamanho e fonte pré-determinadas 
deste modelo. 

\section{Equações, Tabelas, Figuras e Referências}

As equações serão numeradas sequencialmente no texto, com a numeração à direita para identificá-las sempre que for necessário  mencioná-las no texto. Por exemplo, a equação do calor é da da pela Eq.~\ref{Calor} a seguir.   
%
\begin{equation}
\frac{\partial u}{\partial t}-\Delta u = f, \quad \quad  \mathrm{em} \quad \Omega. \label{Calor}
\end{equation}

Os autores podem introduzir tabelas e figuras em seus textos, desde sejam imprescindíveis ao texto e desde que o resumo expandido exceda três páginas, incluindo as referências. As tabelas e figuras devem ter legendas e devem ser mencionadas no texto. A Fig.~\ref{figura01} apresenta um exemplo de inclusão de figuras. utilizando-se a seguinte estrutura:
%
\begin{figure}[ht]
\centering
\includegraphics[width=.1\textwidth]{logo_enset}
\caption{ {\small Logo do ENSET.}}
\label{figura01}
\end{figure}

Para a confecção das tabelas, observe o exemplo da Tabela \ref{tabela01}, conforme segue.
%
\begin{table}[ht]
\caption{ {\small Dados coletados.}}
\centering
\begin{tabular}{ccc}
    \begin{tabular}{c|r|r} % <-- Changed to S here.
      \textbf{Amostra} & \textbf{Valor 1} & \textbf{Valor 2}\\
      \hline
      1 & 1110,1 & 457,9\\
      2 & 10,1 & 7876,2\\
      3 & 23,11 & 896,1\\
    \end{tabular}
\end{tabular}\label{tabela01}
\end{table}

As referências devem estar em ordem alfabética pelo sobrenome do primeiro autor e dos demais, se necessário, usando-se, ainda, ordem cronológica, para trabalhos de um mesmo autor. Trabalhos dos mesmos autores, publicados no mesmo ano, devem ser listados utilizando-se a ordem alfabética do título do trabalho.  Basicamente, as referências devem conter  as iniciais dos nomes dos autores, sendo escrito, por extenso, apenas o último sobrenome. Em seguida, o título do trabalho e o título da publicação (revista, livro, dissertação, tese, anais de evento), volume, páginas,  ano e DOI (se for o caso) ou ISSN (se for o caso). 

A seguir, seguem instruções para diversos tipos de publicações:
\begin{itemize}
\item Quando forem citadas mais de três referências, em sequência, deve-se usar, 
por exemplo, \cite{Diniz1,Diniz2,Jafelice,Santos};
\item Se artigo, após o título da publicação ({\it em itálico}) deve vir o volume e as páginas correspondentes, 
seguidos do ano, conforme as referências \cite{Diniz2, Jafelice}. Trabalhos aceitos, mas não publicados, 
devem ser citados conforme mostrado na referência \cite{Cuminato}. Não havendo DOI, coloque-se (to appear);
\item Se capítulo de livro, após o título da publicação,  deve vir o título da série 
(quando aplicável), o número do capítulo e o volume, como na referência \cite{daSilva};
\item Dissertações, teses e similares devem seguir o padrão das referências \cite{Diniz1,Mallet};
\item No caso de livros, deve-se seguir o padrão da referência \cite{Gomes}, para aqueles publicados 
dentro de uma série, ou \cite{Boldrini};
\item Trabalhos publicados em anais de eventos devem seguir o padrão da referência em \cite{Santos}. 
\end{itemize}

%%%%%%%%%%%%%%%%%%%%%%%%%%%%%%%%%%%%%%%%

\section{Teoremas, Corolários e Outros Ambientes Úteis}\label{sec4}

Teoremas, lemas, proposições, corolários, definições e observações serão automaticamente numerados por seção, conforme os exemplos dados. 

\begin{teorema}\label{sec3_teo1}
Primeiro teorema desta seção.
\end{teorema}

\begin{prop}\label{sec3_prop1}
Primeira proposição desta seção.
\end{prop}

\begin{defi}\label{sec3_defi1}
Primeira definição desta seção.
\end{defi}

\begin{obs}\label{sec3_obs1}
Primeira observação desta seção. As mesmas regras aplicam-se para a 
introdução dos primeiros Lema, Corolário e Exemplo. No modelo em LaTeX, veja o arquivo enset.cls.
\end{obs}


\section{Conclusões}
Resumir, em linhas gerais, as principais conclusões do trabalho. 

\section*{Agradecimentos (opcional)}
Apresentar os agradecimentos às pessoas e instituições e agências de fomento pertinentes.


\begin{thebibliography}{00}

\bibitem{Boldrini} 
J. L. Boldrini, S. I. R. Costa, V. R. Ribeiro, and H. G. Wetzler. {\it Álgebra Linear e Aplicações, 3a. edição}. Harbra, São Paulo, 1984.

\bibitem{Cuminato}
J. A. Cuminato and V. Ruas. Unification of distance inequalities for linear variational problems, 
{\it Comp. Appl. Math.}, 2014. DOI: 10.1007/s40314-014-0163-6.
 
\bibitem{daSilva} 
P. L. da Silva and I. L. Freire. On the group analysis of a modified Novikov equation, 
{\it Interdisciplinary Topics in Applied Mathematics, Modeling and Computational Science}, 
Springer Proceedings in Mathematics and Statistics,  volume 117, chapter 23, pages 161-166, 2015.

\bibitem{Diniz1}
G. L. Diniz. A mudança no habitat de populações de peixes: de rio a represa -- o modelo 
matemático, Dissertação de Mestrado, Unicamp, 1994.

\bibitem{Diniz2}
G. L. Diniz, J. F. C. A. Meyer e L. C. Barros. Solução numérica para um problema de 
Cauchy Fuzzy que modela o decaimento radioativo, {\it TEMA},  23:63--72, 2001. DOI:10.1007/s40314-014-0163-6.

\bibitem{Gomes}
L. T. Gomes, L. C. de Barros, and B. Bede. Fuzzy differential equation in various approaches. 
In {\it SpringerBriefs in Mathematics}. SBMAC- Springer, 2015. ISSN: 2191-8198.

\bibitem{Jafelice} R. M. Jafelice, L. C. Barros and R. C. Bassanezi. Study of the dynamics of HIV under 
treatment considering fuzzy delay, {\it Comp. Appl. Math.}, 33:45--61, 2014.

\bibitem{Mallet}
S. M. Mallet. Análise Numérica de Elementos Finitos. Tese de Doutorado, LNCC/MCTI, 1990.

\bibitem{Santos} I. L. D. Santos e G. N. Silva. Uma classe de problemas de controle ótimo 
em escalas temporais, {\it Proceeding Series of the Brazilian Society of Computational and 
Applied Mathematics}, volume 1, 2013. DOI: 10.5540/03.2013.001.01.0177.

\end{thebibliography}

\end{document}
